\section*{Chapitre \ref{ch:b2:lab-techniques-2} --- Techniques de laboratoire II~: synthèse et analyse en pratique}

\begin{solution}{exo:b2:lab-techniques-2:1}
Eau et glace~; glace et sel~; carboglace et acétone (comme pour les \omterm{def:b2:enolates-aldol:enolate}{énolates} du \cref{ch:b2:enolates-aldol}).
\end{solution}

\begin{solution}{exo:b2:lab-techniques-2:2}
La phase inférieure. Ajouter quelques gouttes d'eau~: elles rejoignent la phase aqueuse, ici la phase supérieure.
\end{solution}

\begin{solution}{exo:b2:lab-techniques-2:3}
Les premières portions s'agglomèrent en captant l'eau~; quand le solide nouvellement ajouté reste libre et
tourbillonne facilement, la solution est sèche.
\end{solution}

\begin{solution}{exo:b2:lab-techniques-2:4}
Le $R_f$ est trop élevé pour une colonne (on veut environ 0.3) et les deux taches sont proches~: utiliser un
éluant moins polaire, plus riche en hexane (9\,:\,1, par exemple), qui abaisse les deux $R_f$ et accroît leur
séparation en volumes de colonne.
\end{solution}

\begin{solution}{exo:b2:lab-techniques-2:5}
\begin{itemize}
  \item Hydrogénocarbonate de sodium (pH 8.3 $>$ 4.2 + 2)~: benzoate dans l'eau~; acidifier, extraire l'acide benzoïque.
  \item Hydroxyde de sodium (pH $>$ 12 $>$ 9.99 + 2)~: phénolate dans l'eau~; acidifier, extraire le phénol.
  \item Acide chlorhydrique (pH voisin de 1 $<$ 4.6 $-$ 2)~: anilinium dans l'eau~; basifier, extraire l'aniline.
  \item Le naphtalène reste dans l'éther~: sécher, évaporer.
\end{itemize}
Le lavage à l'hydrogénocarbonate doit venir en premier~: l'hydroxyde extrairait ensemble l'acide et le phénol.
% ledger: ph:NaHCO3, pka:benzoic, pka:phenol, pka:anilinium
\end{solution}

\begin{solution}{exo:b2:lab-techniques-2:6}
$1/T = 1/T_b - (R/\Delta_{\mathrm{vap}}H)\ln(p/p^\circ)$~: bromobenzène vers \qty{49}{\degreeCelsius},
benzaldéhyde vers \qty{68}{\degreeCelsius} (figure de la \cref{prop:b2:lab-techniques-2:reduced-pressure}).
\end{solution}

\begin{solution}{exo:b2:lab-techniques-2:7}
$0.050 \times 200\,000/400 = \qty{25}{K}$.
\end{solution}

\begin{solution}{exo:b2:lab-techniques-2:8}
Chaque composé a sa propre réponse~: le pourcentage d'aire n'est un pourcentage massique que si les \omterm{def:b2:chromatography:quantitative}{facteurs
de réponse} sont égaux. RMN~: $0.12/3 = 0.040$~mol d'impureté par mol de produit, soit $0.040 \times 100 = \qty{4}{g}$
pour \qty{200}{g} de produit~; $w = 200/204 \approx 98.0~\%$.
\end{solution}

\begin{solution}{exo:b2:lab-techniques-2:9}
Masse théorique~: $0.0300 \times 164 = \qty{4.92}{g}$. Non corrigé~: $4.10/4.92 \approx 83.3~\%$~; corrigé~:
$0.950 \times 4.10/4.92 \approx 79.2~\%$.
\end{solution}

\begin{solution}{exo:b2:lab-techniques-2:10}
La présence de benzène signifie que le réactif s'est formé puis a été protoné~: eau dans la verrerie, le
solvant ou l'aldéhyde~; air (humide) entré par une fuite~; aldéhyde en partie oxydé en acide benzoïque, dont
l'hydrogène acide détruit le réactif. Remèdes~: verrerie séchée à l'étuve, solvant anhydre, aldéhyde
fraîchement distillé, surpression de diazote contrôlée au barboteur. Si presque aucun réactif ne s'était
formé, ce serait l'indice d'un défaut d'amorçage (surface du magnésium couverte d'oxyde), évité par des
copeaux frais ou activés.
\end{solution}

\begin{solution}{exo:b2:lab-techniques-2:11}
\qty{270}{kJ} à \qty{150}{W}~: au moins $270\,000/150 = \qty{1800}{s}$, une demi-heure, et seulement si
l'halogénure réagit aussi vite qu'il est ajouté. Plus vite, le dégagement de chaleur dépasse le
refroidissement~: la température monte, ou du réactif s'accumule et risque de réagir plus tard d'un seul
coup (\cref{prop:b2:lab-techniques-2:accumulation}).
\end{solution}

\begin{solution}{exo:b2:lab-techniques-2:12}
Recristallisation~: $13.0 \times 0.80 = \qty{10.4}{g}$ à 99~\%. Colonne~: $13.0 \times 0.95 =
\qty{12.4}{g}$ à 99.5~\%, mais \qty{1.5}{L} de solvant à acheter, évaporer et éliminer. La colonne donne
environ \qty{2}{g} de plus d'un produit un peu plus pur~; la recristallisation est plus simple, moins chère
et moins génératrice de déchets, et on la préfère quand sa pureté suffit.
\end{solution}

\begin{solution}{pb:b2:lab-techniques-2:1}
\textbf{1.} Éthoxyéthane~: extrêmement inflammable, nocif, forme des peroxydes. Bromobenzène~: inflammable,
toxique, toxique pour les organismes aquatiques. Magnésium~: solide inflammable, libère du dihydrogène au
contact de l'eau.
\textbf{2.} \ce{C6H5MgBr + H2O -> C6H6 + Mg(OH)Br}~: chaque goutte d'eau détruit du réactif~; le séchage à
l'étuve élimine le film d'eau présent sur le verre.
\textbf{3.} Pour tenir à l'écart l'humidité et le dioxygène, qui détruit lui aussi le réactif.
\textbf{4.} Le débit de diazote~; il empêche l'air d'entrer par la sortie.
\textbf{5.} Il ne réagit pas avec le réactif, les doublets de son oxygène stabilisent le magnésium, et sa
basse température d'ébullition (\qty{307.7}{K}) permet au reflux de limiter la température.
\textbf{6.} $15.7/157.01 = \qty{0.1000}{mol}$ de bromobenzène~; $2.67/24.305 = \qty{0.1099}{mol}$ de magnésium,
en excès de 10~\%.
\textbf{7.} $0.1000 \times 270 = \qty{27.0}{kJ}$.
\textbf{8.} Pour que la chaleur soit libérée au rythme auquel le reflux et le bain l'évacuent.
\textbf{9.} $0.0200 \times 270 = \qty{5.40}{kJ}$.
\textbf{10.} $5400/170 \approx \qty{32}{K}$.
\textbf{11.} Environ \qty{52}{\degreeCelsius}, bien au-dessus de la température d'ébullition de l'éthoxyéthane
(\qty{34.6}{\degreeCelsius})~: une ébullition violente qui peut déborder le réfrigérant et projeter des vapeurs
inflammables hors du ballon.
\textbf{12.} Ajouter environ un dixième du bromobenzène, attendre que la réaction démarre (trouble, léger
reflux), puis ajouter le reste au rythme d'un reflux doux, le bain de glace étant prêt.
\textbf{13.} \ce{C6H5MgBr + C6H5CHO} donne l'alcoolate \ce{(C6H5)2CHOMgBr}~; le benzaldéhyde,
$9.55/106.12 = \qty{0.0900}{mol}$, est limitant.
\textbf{14.} L'alcool est doublement benzylique~: un acide fort le transformerait en un cation stabilisé,
conduisant à des éthers ou à des produits de substitution~; le chlorure d'ammonium est assez acide pour
protoner l'alcoolate et dissoudre les sels de magnésium.
\textbf{15.} La phase supérieure~: l'éthoxyéthane est moins dense que l'eau.
\textbf{16.} Il retire l'essentiel de l'eau dissoute dans l'éther.
\textbf{17.} Par du sulfate de magnésium jusqu'à ce qu'il reste pulvérulent, une filtration, puis un
évaporateur rotatif sous pression réduite.
\textbf{18.} Le bromure de phénylmagnésium se couple avec du bromobenzène qui n'a pas réagi~; une forte
concentration locale de bromobenzène et une température élevée le favorisent, autre raison d'un ajout lent.
\textbf{19.} Le biphényle ($R_f$ 0.85), le moins polaire.
\textbf{20.} Environ $1/0.85 \approx 1.2$ volume de colonne pour le biphényle, $1/0.25 = 4$ pour l'alcool.
\textbf{21.} L'alcool apporte 10 protons \omterm{def:b2:huckel:aromatic}{aromatiques} par CH~; l'excédent $10.90 - 10.00 = 0.90$ vient du
biphényle, à 10 protons par molécule~: 0.090 mol de biphényle par mol d'alcool.
\textbf{22.} $w = 184.24/(184.24 + 0.090 \times 154.21) \approx 0.930$.
\textbf{23.} $12.50 \times 0.930 \approx \qty{11.62}{g}$.
\textbf{24.} $0.0900 \times 184.24 \approx \qty{16.58}{g}$.
\textbf{25.} $12.50/16.58 \approx 75.4~\%$.
\textbf{26.} $11.62/16.58 \approx \boldsymbol{70~\%}$ (70.1~\%), corrigé de la pureté.
\end{solution}
